% =============================================================================
%  preambulo.tex — Preámbulo de los resúmenes en español
%  Evsukoff, A. (2020). Inteligência Computacional. Rio de Janeiro: E-papers.
% =============================================================================
\documentclass[11pt,a4paper]{article}

\usepackage[T1]{fontenc}
\usepackage[utf8]{inputenc}
\usepackage[spanish,es-tabla]{babel}
\usepackage{lmodern}
\usepackage[a4paper,margin=2.4cm]{geometry}
\usepackage{amsmath,amssymb,amsthm,mathtools}
\usepackage{bm}
\usepackage{enumitem}
\usepackage{xcolor}
\usepackage{tcolorbox}
\tcbuselibrary{skins,breakable}
\usepackage{hyperref}
\usepackage{microtype}
\usepackage{titlesec}
\usepackage{fancyhdr}
\usepackage{booktabs}
\usepackage{array}

\hypersetup{
  colorlinks=true,
  linkcolor=black!70,
  urlcolor=blue!60!black,
  citecolor=black!70,
  pdfborder={0 0 0}
}

\titleformat{\section}{\Large\bfseries\sffamily\color{black!85}}{\thesection}{0.7em}{}
\titleformat{\subsection}{\large\bfseries\sffamily\color{black!75}}{\thesubsection}{0.6em}{}

\pagestyle{fancy}
\fancyhf{}
\fancyhead[L]{\small\itshape Resumen — Inteligencia Computacional (Evsukoff, 2020)}
\fancyhead[R]{\small\thepage}
\renewcommand{\headrulewidth}{0.3pt}

\newtcolorbox{conceito}[1][]{
  colback=blue!4, colframe=blue!55!black, boxrule=0.6pt,
  arc=2mm, left=6pt, right=6pt, top=4pt, bottom=4pt,
  title={\bfseries Concepto clave}, fonttitle=\sffamily\bfseries,
  breakable, #1
}
\newtcolorbox{formula}[1][]{
  colback=gray!6, colframe=gray!50, boxrule=0.4pt,
  arc=1.5mm, left=6pt, right=6pt, top=3pt, bottom=3pt,
  breakable, #1
}
\newtcolorbox{aplicacao}[1][]{
  colback=green!4, colframe=green!45!black, boxrule=0.5pt,
  arc=2mm, left=6pt, right=6pt, top=3pt, bottom=3pt,
  title={\bfseries Aplicación}, fonttitle=\sffamily\bfseries,
  breakable, #1
}

\newcommand{\R}{\mathbb{R}}
\newcommand{\E}{\mathbb{E}}
\newcommand{\Var}{\operatorname{Var}}
\newcommand{\Cov}{\operatorname{Cov}}
\newcommand{\diag}{\operatorname{diag}}
\newcommand{\argmin}{\operatorname*{arg\,min}}
\newcommand{\argmax}{\operatorname*{arg\,max}}
\newcommand{\T}{^{\!\top}}
\DeclareMathOperator{\sign}{sign}
\DeclareMathOperator{\tr}{tr}
